\section{Overall Research Design}

This thesis proposes and evaluates a hybrid approach to automated
FAIRness assessment of research software repositories. The approach
combines two complementary assessment components: a rule-based
component that evaluates structured repository evidence against
indicators derived from the FAIR4RS principles
\cite{barker2022fair4rs}, and an LLM-based component that evaluates
predefined qualitative characteristics of README documentation.
Together, these components provide complementary evidence for
assessing different aspects of software FAIRness.

The motivation for a hybrid design stems from the complementary
strengths and limitations of existing assessment approaches.
Rule-based tools such as F-UJI \cite{fuji2021} and howfairis assess
machine-detectable evidence such as licensing, persistent
identifiers, citation information, and other structured repository
characteristics, but do not assess qualitative README characteristics
such as installation guidance, usage instructions, examples, or
contribution information. The rule-based component therefore provides
deterministic and reproducible assessments, while the LLM-based
component addresses documentation characteristics that are difficult
to capture through fixed presence-based rules.

The research design follows four main phases. First, a dataset of 
188 research software repositories is collected from GitHub using 
a systematic topic-based search strategy. Second, a ground truth 
set of 30 repositories is manually annotated against a structured 
rubric derived from FAIR4RS principles, providing a reference for 
evaluating the accuracy of both assessment components. 
Third, the rule-based and LLM-based components are implemented 
and validated against the ground-truth annotations, with the LLM-based 
component evaluated using two separate large language model configurations.
The rule-based component is subsequently applied to the full 188-repository 
dataset, while the LLM-based evaluation is conducted on the 30-repository 
annotated corpus for both LLM configurations. Fourth, the outputs of both 
components are evaluated to assess agreement between automated assessments and manual 
annotations, and to compare the proposed approach with existing baseline tools.
The following sections describe each component of the methodology 
in detail.

\section{Dataset Collection}

\subsection{Research Software Definition}

Research software is defined in this work as software developed to 
support, enable, or conduct scientific research, including data 
analysis tools, simulation software, computational pipelines, 
scientific frameworks, and domain-specific research libraries 
\cite{barker2022fair4rs}. Repositories were considered research 
software if they demonstrated a clear scientific or research context 
through their descriptions, repository topics, README documentation, 
associated publications, or the affiliation of the owning organisation.

\subsection{Collection Strategy}

To obtain a diverse dataset, repositories were collected from ten 
research domains: bioinformatics, computational biology, astronomy, 
neuroscience, climate science, scientific computing, research 
software, cheminformatics, computational physics, and reproducible 
research. Each domain was queried using the GitHub Search API with 
the following pattern:

\begin{quote}
\texttt{topic:<domain> stars:>20 archived:false}
\end{quote}

A minimum threshold of 20 GitHub stars was applied to filter out 
inactive or personal experimental repositories. Archived repositories 
were excluded as they are no longer actively maintained. The 
collection was capped at 20 repositories per topic to ensure balanced 
domain representation, yielding an initial dataset of 198 repositories.

\subsection{Filtering and Cleaning}

Following collection, repositories were automatically filtered 
according to predefined inclusion criteria. Forks, archived 
repositories, repositories whose most recent push occurred before 
1 January 2023, and repositories whose names or descriptions 
contained predefined keywords --- such as \textit{awesome}, 
\textit{tutorial}, \textit{course}, \textit{template}, 
\textit{workshop}, \textit{lecture}, \textit{homework}, 
\textit{assignment}, and \textit{cheatsheet} --- were excluded, 
as these typically represent curated lists, educational materials, 
or project templates rather than executable research software. 
These filtering steps reduced the initial dataset from 198 to 
188 repositories.

\subsection{Dataset Characteristics}

The final dataset contains 188 research software repositories 
spanning 10 research domains and multiple programming language 
ecosystems. Python is the most common language (86 repositories, 
46\%), followed by C++ (29, 15\%) and Jupyter Notebook (18, 10\%). 
Additional languages include Julia, R, Rust, Java, Fortran, Go, 
and JavaScript, providing diversity beyond a single programming 
ecosystem. All repositories were publicly accessible on GitHub and 
had been actively maintained at the time of collection, whose latest 
repository activity occurred on or after 1 January 2023.

For each repository, the following data were collected via the 
GitHub API: repository metadata (name, owner, description, primary 
language, star count, topics, and last push date), README file 
content, and the presence of FAIR-related repository files including 
license files, citation files, dependency specifications, structured 
metadata files, and contribution guidelines.

\section{Ground Truth Construction}

\subsection{Purpose and Scope}

To evaluate the performance of both the rule-based and LLM-based
assessment components, a manually annotated reference corpus of
30 research software repositories was constructed. Of these,
23 repositories are also included in the final 188-repository
dataset, while seven originate from the earlier annotation corpus
and are not members of the final collected dataset. Manual
annotations provide reference scores against which the automated
assessments can be compared, enabling the computation of agreement
metrics and the identification of systematic scoring errors.

\subsection{Repository Selection}

The annotated corpus was constructed in two stages. An initial set
of 10 repositories was annotated during the pilot phase of system
development and was used to refine the annotation guidelines and
assessment procedure. The corpus was subsequently expanded to
30 repositories to increase diversity across research domains and
programming-language ecosystems, including C++, Java, MATLAB,
Fortran, Julia, Go, Rust, and Groovy in addition to Python.

The final annotated corpus therefore contains both repositories from
the initial pilot and repositories added during the subsequent
expansion. Of the 30 annotated repositories, 23 are also represented
in the final 188-repository dataset. The remaining seven were retained
from the earlier annotation corpus but are not members of the final
188-repository collection.

Because the original 10 pilot repositories contributed to rubric
development and prompt refinement, evaluation on the full
30-repository corpus is not treated as a fully independent test of
the LLM-based component. Chapter~5 therefore additionally reports a
pilot-exclusion sensitivity analysis in which these repositories are
excluded from the LLM agreement calculations, leaving 20 repositories
that were not used during the original pilot phase.
\subsection{Annotation Guidelines}

Annotations were performed using a structured rubric developed 
specifically for this thesis, covering both the rule-based 
indicators (F1, F2, A1, I1, R1, R2) and the AI-based documentation 
quality dimensions (D1--D6). The rubric defines precise scoring 
criteria for each indicator and dimension, including edge case 
handling for non-standard license filenames, alternative citation 
file formats, and the treatment of external documentation links.

The annotation guidelines were developed iteratively during the
pilot phase. Comparison with automated scorer outputs was used to
identify ambiguous cases and inconsistencies in the operational
definitions, which were then resolved against repository evidence and
the intended FAIR4RS interpretation. The final version, v1.4, clarified the distinction between 
structured citation metadata (F2) and explicit citation instructions 
(R2), and resolved ambiguities identified during ground truth 
annotation. All 30 repositories were annotated against the 
finalised v1.4 guidelines to ensure consistency.

\subsection{Annotation Process}

Each repository was annotated by the thesis author by examining 
the repository root file structure and README content. Rule-based 
indicators (F1--R2) were scored by checking for the presence of 
specific files and patterns in the README text. AI-based dimensions 
(D1--D6) were scored on a 0--2 scale by reading the README and 
assessing the clarity and completeness of installation instructions, 
usage guidance, purpose description, dependency information, 
examples, and contribution information.

Each annotation record includes the assigned score together with 
evidence notes justifying the decision. For the AI-based dimensions, 
annotations additionally record the annotator's confidence level 
(HIGH, MEDIUM, or LOW), reflecting the degree of ambiguity 
encountered when applying the rubric criteria.

\subsection{Inter-Annotator Agreement}

To assess annotation reliability, a second annotator was provided 
with the finalised v1.4 annotation guidelines, the list of 30 
repositories with GitHub links, pre-fetched README files, 
FAIR-related file presence summaries, and blank annotation forms. 
The annotation package was designed to enable fully independent 
annotation without requiring access to the primary annotations or 
the GitHub repositories directly.

Upon completion of the independent annotation process,
inter-annotator agreement was quantified using Cohen's weighted
Kappa for the AI-based dimensions (D1--D6), where partial agreement
is meaningful, and Cohen's unweighted Kappa for the binary
rule-based indicators (F1--R2). Exact agreement was additionally
reported to aid interpretation in cases where Kappa was undefined
because of a lack of class variation.

Disagreements and anomalous results were subsequently reviewed
against the final annotation guidelines and the collected repository
evidence. This validation identified, among other cases, an
inconsistency between the file-detection schema originally used for
the ground-truth collection and the expanded schema used for the
full dataset. The ground-truth file inventories and affected primary
annotations were corrected where the repository evidence
unambiguously supported the correction. The final agreement results
are reported in Chapter~5.

\section{Rule-Based Assessment}

\subsection{Overview}

The rule-based assessment component evaluates each repository against 
six binary indicators derived from the FAIR4RS principles 
\cite{barker2022fair4rs} and the FAIR-IMPACT research software 
metrics \cite{chuehong2025frsm}. Each indicator checks for the 
presence of repository files or textual patterns that provide 
evidence for selected FAIR4RS indicators. The component produces a 
score of 0 or 1 per indicator, with a maximum total score of 6.

\subsection{Indicators}

The six rule-based indicators are defined as follows:

\begin{description}
\item[F1 --- Persistent Identifier] Score 1 if evidence of a persistent identifier (e.g., DOI, Zenodo,
or Software Heritage identifier) is detected in the repository
documentation or metadata. Detected patterns include 
doi.org URLs, Zenodo badges, Software Heritage identifiers, and 
DOI badge URL formats.

\item[F2 --- Structured Metadata] Score 1 if a structured metadata 
file is present in the repository root. Accepted metadata files are 
\texttt{CITATION.cff} and \texttt{codemeta.json}, both of which 
provide machine-readable software metadata.

\item[A1 --- License] Score 1 if licensing information is explicitly
declared for the repository. This condition is satisfied if a
recognized license file is present in the repository root, if a
license is declared in the GitHub repository metadata, or if the
README explicitly states the applicable license. Accepted license
filenames include \texttt{LICENSE}, \texttt{LICENSE.md},
\texttt{LICENSE.txt}, \texttt{LICENSE.rst}, \texttt{COPYING}, and
\texttt{COPYING.txt}.


\item[I1 --- Dependency Specification] Score 1 if a dependency 
specification file is present in the repository root. Accepted 
files cover multiple language ecosystems, including 
\texttt{requirements.txt}, \texttt{pyproject.toml}, 
\texttt{environment.yml}, \texttt{setup.py}, \texttt{Cargo.toml}, 
\texttt{go.mod}, \texttt{pom.xml}, \texttt{Project.toml}, and 
\texttt{package.json}.

\item[R1 --- README Present] Score 1 if a README file is present 
in the repository root. Accepted filenames include README.md, 
README.rst, README.txt, and README.

\item[R2 --- Citation Information] Score 1 if explicit citation 
instructions are present in the repository. This includes dedicated 
citation files (CITATION.md, CITATION.bib, CITATION.txt, CITATION, 
CITATION.cff) and explicit citation instructions in the README text, 
such as phrases containing ``please cite'', ``how to cite'', ``if you 
use \ldots{} cite'', or ``we would appreciate citations''.
\end{description}

\subsection{Implementation}

Repository data were collected in a preprocessing step and stored 
locally as JSON files. The scorer operates exclusively on these 
locally stored artefacts rather than querying the GitHub API during 
scoring, ensuring deterministic and reproducible execution. For each 
repository, the scorer reads the FAIR file presence record 
(\texttt{fair\_file\_checks.json}), the repository metadata 
(\texttt{metadata.json}), and the README file. File presence checks are performed against the stored JSON record,
while selected indicators additionally inspect repository metadata
and textual evidence. In particular, F1 searches for persistent
identifier patterns, A1 checks repository metadata and explicit
license statements in the README, and R2 searches for explicit
citation instructions in the README text.

The scorer was validated against the 30 manually annotated ground
truth repositories. Initial validation revealed discrepancies in
several indicators. Detection logic for F1 and R2 was refined to
cover additional DOI representations, citation file formats, and
citation instruction phrasings identified during manual validation.

Accordingly, the rule-based validation against this corpus is
interpreted primarily as a consistency check between the finalised
rubric, repository evidence, and deterministic implementation rather
than as a fully independent held-out evaluation.

A subsequent consistency check also identified schema drift between
the collector originally used for the ground-truth repositories and
the collector used for the full dataset. The earlier collector did
not check several language-specific dependency and build files,
including \texttt{Cargo.toml}, \texttt{go.mod}, \texttt{pom.xml},
\texttt{Project.toml}, and \texttt{build.gradle}. The ground-truth
file inventories were therefore aligned with the expanded schema
used for the full dataset, and affected annotations were corrected
where the repository evidence unambiguously satisfied the final
rubric. Additional manual annotation inconsistencies identified
during validation were likewise corrected against the finalised
guidelines. The final validation results, including precision,
recall, and F1-score per indicator, are reported in Chapter~5.

\subsection{Output}

The rule-based scorer produces a CSV file containing one row per 
repository with individual indicator scores and the aggregated 
rule-based score. The scorer was subsequently applied to all 188
repositories in the dataset, producing the rule-based assessment 
results used in the evaluation presented in Chapter~5.

\section{LLM-Based Assessment}

\subsection{Overview}

The LLM-based assessment component evaluates the quality and 
completeness of repository README documentation against six 
documentation quality dimensions derived from FAIR4RS reusability 
indicators \cite{barker2022fair4rs} and the software documentation 
quality framework proposed by Treude et al. \cite{treude2020beyond}. 
Unlike the rule-based component, which checks for the presence of 
files and patterns, the LLM-based component assesses the content 
and clarity of documentation that a potential user would encounter 
when attempting to reuse the software.

\subsection{Documentation Assessment Dimensions}

Six qualitative documentation dimensions are evaluated, each scored
on a 0--2 ordinal scale:

\begin{description}
\item[D1 --- Installation Clarity] Assesses whether a new user can
install the software from the README alone. Score 2 if complete and
actionable installation instructions, including relevant
prerequisites, are provided; score 1 if installation information is
present but incomplete or relies on external documentation; score 0
if no installation instructions are present.

\item[D2 --- Usage Clarity] Assesses whether a new user can 
understand how to run or use the software from the README. Score 
2 if a concrete command, code snippet, or executable example is 
present; score 1 if external documentation is clearly linked but 
no inline example appears; score 0 if no usage information is 
present.

\item[D3 --- Purpose Description] Assesses whether the README 
clearly explains what the software does and its intended domain or 
audience. Score 2 if the purpose is clearly described; score 1 if 
a description is present but vague or incomplete; score 0 if no 
description is present.

\item[D4 --- Dependency Description] Assesses whether the README 
describes the environment or dependencies required to use the 
software. Score 2 if dependencies are explicitly listed in 
human-readable form; score 1 if dependency information is only 
partially provided; score 0 if no dependency information is present.

\item[D5 --- Examples and Tutorials] Assesses whether the README
provides concrete examples that demonstrate software use. Score 2 if
a meaningful usage example is provided directly in the README.
Score 1 if the README explicitly references example-oriented
resources such as tutorials, galleries, cookbooks, or notebooks, but
their content is not included in the supplied evidence. Generic
documentation, installation, testing, contribution, or development
links alone do not count as examples. Score 0 if no qualifying
examples or tutorials are present.

\item[D6 --- Contribution Information] Assesses whether the README 
explains how to contribute to or get support for the software. 
Score 2 if contribution guidance is provided together with at 
least one community or support channel; score 1 if contribution 
or support information is only partially present; score 0 if 
neither is available.
\end{description}

\subsection{Model Selection}

The open-source Llama 3.1 8B model was selected because it can be 
executed locally through Ollama, ensuring reproducibility and 
eliminating dependence on commercial APIs. Running the model locally 
also avoids API usage costs and enables repeated experiments under 
identical inference settings. Temperature was set to 0 to reduce sampling variability and promote
more consistent outputs across runs.
To examine whether findings observed with the primary Llama
configuration were specific to that model, a second model,
\texttt{Qwen3.6-35B-A3B-MLX-8bit}, from the Qwen model family
\cite{yang2025qwen3}, was evaluated using the same substantive
rubric, README evidence policy, and ground-truth corpus. The exact
identifier refers to the model configuration exposed by the
university KISTE service and is reported for reproducibility; the
cited Qwen family publication does not document this exact
KISTE-hosted configuration.

The model was accessed through an OpenAI-compatible API endpoint
hosted on a departmental GPU server rather than executed locally.
Recurring patterns across both models are interpreted as evidence
that the corresponding finding is not unique to Llama~3.1~8B,
whereas divergent results indicate sensitivity to model choice.
The comparison is limited to two model configurations and is not
intended to establish model-independent properties of LLM-based
FAIRness assessment.

\subsection{Prompt Design}

The LLM is provided with the repository name and README content as 
input and instructed to score each dimension according to the 
defined rubric. The prompt specifies that the model must evaluate 
only evidence explicitly present in the README text, and must not 
assume content exists behind external links. Each score must be 
accompanied by a confidence level (HIGH, MEDIUM, or LOW) and a 
one-sentence reasoning statement citing specific README evidence.

The final prompt was developed through iterative experimentation on 
the pilot repository set, with refinements improving the specificity 
of dimension definitions, particularly for D5, where generic 
documentation links were being incorrectly treated as usage examples.

For the Qwen evaluation, the same substantive dimension definitions, 
scoring criteria, and README-only evidence restrictions were retained. 
Backend-specific output instructions were used from the outset to obtain 
the required JSON structure; no change to the rubric or scoring criteria 
was made when adapting the prompt to this model.

\subsection{Implementation}

For the primary Llama evaluation, the LLM-based scorer uses the
Ollama HTTP API with structured JSON output enforced through a JSON
schema supplied to the API, ensuring that every response conforms to
the expected structure and that scores are restricted to the
permitted values (0, 1, or 2).
README content is supplied directly to the model as assessment
evidence. To bound the amount of input passed to the model and keep
preprocessing consistent across repositories, the implementation
applies a fixed upper limit of 12,000 characters. For READMEs
exceeding this threshold, only the first 12,000 characters are
supplied to the model.

The threshold is a practical implementation-level preprocessing
choice rather than a limit derived from the model's maximum context
capacity or a threshold selected through experimental optimisation.
Character count was used because it provides a simple and
deterministic way to bound input without introducing an additional
tokenization step. However, language models process tokens rather
than characters, and the relationship between character count and
token count varies with the content of the README. A token-based
limit would therefore provide a more direct measure of the actual
model input size.

The implementation also applies the limit as a sequential cutoff
rather than selecting content according to README structure. This
approach is reproducible across repositories but may exclude relevant
documentation when important sections occur after the cutoff. A
section-aware strategy that identifies and preserves relevant
sections, such as installation, usage, dependencies, examples, and
contribution guidance, would be more closely aligned with the D1--D6
assessment dimensions.

Within the 30-repository ground-truth corpus, only one README
(\texttt{Ouranosinc/xclim}) exceeded the 12,000-character threshold,
with a total length of 14,262 characters. Inspection of the excluded
portion showed that it consisted of reStructuredText definitions for
images, badges, logos, and associated links rather than substantive
documentation corresponding to D1--D6. No relevant assessment
evidence was therefore identified as lost through truncation in this
case. Nevertheless, this observation concerns a single repository
and does not establish that fixed character-based truncation is
generally safe for longer READMEs. Token-aware and section-aware
input handling are therefore more principled alternatives for future
versions of the framework.

For Qwen3.6-35B-A3B-MLX-8bit, the same rubric and dimension definitions 
were used, but the scorer was implemented against an OpenAI-compatible chat 
completion endpoint rather than the Ollama HTTP API, since the model was hosted 
remotely rather than executed locally. Unlike the Ollama API, this endpoint offered 
no verified mechanism for enforcing a JSON schema on the response, so structured output 
relied on explicit prompt instructions combined with post-hoc validation of the returned 
JSON, including defensive stripping of Markdown code-fence wrappers that this model occasionally 
produced. This model additionally returns an internal reasoning trace as a separate field from its 
final answer; only the final answer content was used for scoring, and the reasoning trace was not 
incorporated into the D1--D6 dimensions.

An initial output-length limit of 4{,}096 tokens resulted in invalid, incomplete, or unparsable 
outputs for seven of the 30 repositories. These seven requests were retried with an increased 
output limit of 16{,}000 tokens and an additional instruction constraining each per-dimension 
reasoning statement to a short sentence. The README input, D1--D6 rubric, score definitions, 
evidence policy, and temperature remained unchanged. The 23 previously successful outputs 
were retained, and the recovery run produced valid outputs for the seven remaining repositories, 
resulting in a complete 30-repository Qwen result set.

Confidence labels and reasoning statements were retained for
qualitative error analysis. Confidence labels were not incorporated
into the D1--D6 documentation scores themselves; instead, their
relationship with assessment correctness was evaluated separately
through the confidence analysis described below.

\subsection{External Link Policy}

The LLM receives only the README text as input and does not retrieve
external documentation linked from the README. Consequently, linked
resources whose content is not included in the supplied evidence are
not treated as equivalent to documentation directly available in the
README. The scoring rubric accounts for such references explicitly;
for example, a clearly linked external resource may provide partial
evidence for dimensions where the referenced content itself is not
available to the model.

Human annotators could account for references to such resources
according to the annotation procedure, whereas the automated
assessment was restricted to the README text supplied to the model.
This evidence-boundary difference is therefore considered when
interpreting disagreements between human and automated assessments.

\subsection{Output}

The LLM scorer produces a CSV file per evaluated model, containing 
one row per repository with scores, confidence levels, and reasoning 
for each of the six dimensions, together with the model name and a 
timestamp. Both scorers were applied to all 30 ground truth repositories, 
producing the two LLM-based assessment result sets used in the evaluation 
presented in Chapter~5.

\section{Repository FAIRness Profile}

The hybrid assessment framework reports the outputs of the
rule-based and LLM-based components together as a unified FAIRness
profile for each repository. This representation maintains the
interpretability of both structural FAIR indicators and qualitative
documentation characteristics, allowing users to identify strengths
and weaknesses across different aspects of software FAIRness.

The rule-based component contributes six binary indicator scores 
(F1--R2) covering structural metadata compliance, together with 
an aggregated rule-based total. The LLM-based component contributes 
six documentation quality scores (D1--D6) on a 0--2 ordinal scale, 
each accompanied by a confidence level and a reasoning statement.

Rather than aggregating these heterogeneous structural and semantic
indicators into a single scalar score, the framework preserves them
as separate dimensions. This avoids introducing arbitrary weights
across dimensions with fundamentally different evidence types ---
binary file-presence checks versus qualitative, LLM-derived
documentation judgments --- and prevents this incomparable evidence
from being collapsed into a potentially misleading overall FAIRness
value. Both components are stored together within a structured assessment
record for each repository. The generated results are persisted as
CSV output for subsequent analysis and are additionally exposed
through the web-based prototype described in Chapter~4. The interface
presents the structural and documentation components separately and
allows the underlying documentation assessment reasoning, confidence
labels, scoring criteria, and available README evidence to be
inspected.

This parallel representation provides more actionable feedback than 
a single aggregated score would allow. A repository may score 
highly on structural metadata indicators while exhibiting weak 
documentation quality, or conversely may have well-written 
documentation but lack persistent identifiers or structured 
citation metadata. Presenting both dimensions separately makes 
these patterns visible and directly supports the identification 
of specific FAIRness gaps.

The resulting hybrid FAIRness profiles constitute the primary
output of the proposed framework, providing a unified representation
of structural and documentation-based assessment results that is
used throughout the evaluation presented in Chapter~5.

In addition to presenting the saved assessments used in the empirical
study, the implemented prototype can apply the assessment procedure
to a newly supplied public GitHub repository. For live documentation
assessment, the prototype supports Llama~3.1~8B through a local
Ollama instance and \texttt{Qwen3.6-35B-A3B-MLX-8bit} through the
university KISTE service. Such live assessments use current repository
content and are kept separate from the fixed datasets and evaluation
results used in this thesis.

\section{Evaluation Strategy}

The evaluation strategy assesses the performance, reliability, and
complementary contribution of the proposed hybrid assessment
framework through quantitative and qualitative analyses. The evaluation is
structured around the two research questions defined in
Chapter~1 and is conducted against the 30-repository ground-truth
corpus described in Section~3.3.

Where applicable, this evaluation is conducted separately for each 
of the two evaluated LLM configurations, Llama~3.1~8B and Qwen3.6-35B-A3B-MLX-8bit. 
For both configurations, agreement is additionally recomputed after excluding 
the ten repositories used during pilot development; this 20-repository analysis 
is treated as a pilot-exclusion sensitivity analysis rather than an independent 
held-out evaluation.

\subsection{Ground Truth Comparison}

The primary evaluation compares the automated scores produced by
both the rule-based and LLM-based components against the manual
ground truth annotations. For the rule-based indicators (F1--R2),
which are binary, performance against the manual ground truth is
evaluated using precision, recall, and F1-score. Exact agreement and
Cohen's unweighted Kappa are additionally reported to characterise
agreement between assessments. For the LLM-based dimensions
(D1--D6), which use a 0--2 ordinal scale, agreement is measured
using exact match accuracy, Mean Absolute Error (MAE), and Cohen's
weighted Kappa per dimension, where partial agreement is given
proportional credit.

\subsection{Inter-Annotator Agreement}

To assess the reliability of the manual ground truth annotations,
inter-annotator agreement between the primary annotator and the
second annotator was quantified using Cohen's unweighted Kappa for
binary indicators and Cohen's weighted Kappa for ordinal dimensions.
Exact agreement was also reported, particularly because Kappa is
undefined when both annotators assign only a single class. These
results establish the reliability of the manual reference
annotations and contextualise the agreement observed between the
automated assessments and the ground truth.

\subsection{Baseline Comparison}

The proposed approach is compared with existing FAIR assessment
tools. A direct comparison with howfairis is performed on the
ground truth repositories, examining indicator coverage, agreements
and disagreements, and FAIRness aspects not represented by the
existing tool. F-UJI is separately evaluated in a pilot experiment
by submitting plain GitHub repository URLs directly, to investigate
whether a dataset-oriented FAIR assessment tool can serve as a
meaningful quantitative baseline when applied to research software
repositories without a mediating dataset record.

\subsection{Qualitative Error Analysis}

In addition to quantitative agreement metrics, representative 
cases of agreement and disagreement between the automated 
assessments and the manual annotations are analysed. The 
confidence levels and reasoning statements produced by the LLM 
are used to help interpret discrepancies and identify common 
sources of scoring differences.

\subsection{Confidence Analysis}

The usefulness of each LLM's self-reported confidence labels is
evaluated separately by comparing prediction correctness across HIGH,
MEDIUM, and LOW confidence outputs. Agreement with the primary human
annotations is examined both across all documentation dimensions
and within individual dimensions where sufficient observations are
available. This analysis assesses whether confidence can be
interpreted as an indicator of prediction reliability or should
instead be treated only as contextual information.

\subsection{README Length and Verbosity Bias}

To investigate whether LLM scoring errors are associated with README
length, the analysis is performed separately for each evaluated LLM
configuration. For each model, the difference between the LLM and
human documentation scores is compared with the effective README
length supplied to the model.
README length is measured after applying the same
12,000-character input limit used during inference. Spearman's rank
correlation is computed for the aggregate documentation score and
separately for D1--D6. This analysis tests whether longer README
files are systematically associated with higher LLM scores
independently of the corresponding human assessments.
